% backmatter/epilogue.tex

\addchap{Эпилог: Как отслеживать изменения и не читать устаревшее}\label{ch:change-tracking}
Платёжная отрасль создаёт особый вид усталости: даже после того,
как архитектура понятна и протоколы изучены, команда всё равно остаётся под риском
устаревшего знания. Документ переиздан, порог обновлён, интерфейс
для участников поменялся, бюллетень уточнил старое правило~--- а~во внутренней базе знаний всё ещё лежит прошлогодний пересказ. Из-за этого правильная по замыслу система начинает принимать неверные решения, опираясь на пересказ источника истины вместо самого источника.

Платёжная индустрия~--- подвижная среда, в~ней может измениться что угодно: справочные детали, разделение
ответственности, процедура оспаривания, допустимый сценарий
аутентификации или форма технической интеграции. Отсюда три задачи: быстро выйти на основной источник, отделить устойчивую
концепцию от быстро стареющей справки и выстроить процесс
проверки изменений.

В~платёжной индустрии ошибаются по двум причинам: из-за того, что не нашли нужный документ,
и~из-за того, что не отследили его обновление. Команда открывает нужную спецификацию,
находит нужный раздел и~всё равно принимает неверное решение:
перед ней редакция прошлого квартала. Поэтому работа с~источниками
опирается на устойчивый порядок
обновления знания: кто следит за изменениями, как их фиксирует
и в~какой момент переводит в~инженерную задачу.

\section*{Где меняется}\label{sec:ps-tracking-layers}

Когда в~один и~тот же месяц приходят бюллетень карточной сети,
разъяснение PCI SSC и~письмо регулятора, их легко сложить в~одну
папку с~пометкой \enquote{надо изучить}. Но эти сигналы несут разную опасность, и~сначала нужно понять, в~какой плоскости произошло изменение.

Первая плоскость~--- правила участия: новые обязанности банков,
пороги мониторинговых программ, ограничения на повторы, пересмотр
логики диспутов. Вторая~--- протокол: новые поля, новая
версия 3DS, изменение соответствия полей в~сообщениях, требования
к~токенам или к~ядру терминала. Третья~--- комплаенс: новые
контрольные меры, даты прекращения прежних режимов, уточнение зоны охвата.
Четвёртая плоскость~--- регулятивная: меняется допустимость продукта, вступает в~силу закон, положение
или разъяснение регулятора.

От характера изменения зависит, кто берёт его в~работу, какой горизонт
реакции допустим и~чем грозит промедление. Протокольная правка
проявляется как падение доли одобрений или сбой приёма карт;
изменение в комплаенсе проявляется во время аудита;
регуляторная перемена проявляется как вопрос
о допустимости работы в~таком виде. Первый шаг при чтении
нового документа~--- определить, где сдвинулась модель.

\section*{Какие сигналы читать заранее}\label{sec:ps-tracking-bulletins}

У~главных документов платёжной отрасли обязательность редко появляется внезапно. У Visa это хорошо видно
по Merchant Business News Digest: обновления выходят с~указанием
региона и~даты вступления в~силу. У Mastercard изменения
выходят новыми редакциями Rules и Transaction Processing Rules,
а~внутри документов для отдельных требований прямо фиксируются
региональные даты вступления в
силу~\cite{visa-business-news-digest,mc-rules,mc-tpr}.
Значит, документы сети нужно читать заранее,
в~темпе самих публикаций.

PCI SSC устроен похоже, хотя форма сигналов иная: новые версии
стандарта, требования с~заранее указанной датой обязательности,
дополнительные разъяснения и отраслевые
бюллетени~\cite{pci-dss-4,pci-faq-1585,pci-format4-suspension}.
EMVCo ведёт архив спецификаций и связанных с~ними бюллетеней,
в~котором рядом видны версия, дата публикации и тип
документа~\cite{emvco-bulletins-archive}.
Банк России публикует проекты нормативных актов в~разделе
публичного обсуждения и указывает сроки для замечаний ещё
до принятия документа~\cite{cbr-project-na}.

Самая неприятная ситуация~--- когда публичный материал даёт только сигнал, а~полный текст требования закрыт. Visa прямо
пишет, что этот дайджест состоит из кратких сводок
и~не заменяет Visa Business News, а~сами статьи закрыты
для внешних читателей. Mastercard, в~свою очередь, выносит
материалы, обновления, объявления, FAQ и бюллетени в Chip
Information Center для участников, зарегистрированных
в Connect~\cite{visa-business-news-digest,mc-chip-info-center}.
Пока закрыт основной текст, достраивать реализацию по пересказам~--- ошибка.
Сначала нужно зафиксировать сам факт изменения и дату вступления
в~силу, затем получить основной документ через участника сети,
эквайера или процессингового партнёра и~после этого
проектировать решение.
\importantnote{%
  Если между открытыми и закрытыми материалами остаётся разрыв, команда
  начинает достраивать детали по форумам и пересказам. Это заканчивается
  лишней переделкой на внедрении или сертификации.

}
\section*{Как превращать новое правило в~задачу}\label{sec:ps-tracking-triage}

Новый документ становится полезен, когда команда превратила его в~задачу. Для этого достаточно короткого разбора.

\begin{enumerate}
  \item К~какой плоскости относится изменение: правила участия, протокол,
    комплаенс или регулирование?
  \item Какие роли и~системы оно затрагивает: issuer, acquirer,
    шлюз, функцию управления рисками, сверку, юридическую команду?
  \item Нужна ли правка кода, процесса, документации или пока
    достаточно мониторинга?
  \item Каков крайний срок внедрения и чем грозит его пропуск?
\end{enumerate}

В~этот момент документ перестаёт быть новостью и становится
конкретной задачей в~рабочей очереди. Без такого перевода изменения
либо игнорируют до последнего срока, либо в~спешке переделывают
не ту часть системы.

\practicenote{%
  Полезен короткий формат разбора изменения: источник,
  основной документ, дата вступления в~силу, затронутые системы,
  требуемое действие и~ссылка на внутреннюю задачу. Такая запись
  переживает почту и~чаты лучше длинного конспекта.

}
\section*{Что оставлять в~книге, а~что сверять заново}\label{sec:ps-tracking-book-vs-reference}

Здесь возникает редакционный вопрос: книга не должна претендовать на вечную актуальность там, где отрасль обновляется короткими
циклами. Материал в~ней делится на два класса по сроку жизни: длинный и~короткий.

В~книге должны оставаться архитектурные связи, роли участников,
логика протоколов и устойчивые инварианты. Сведения же, которые
живут один-два цикла обновления правил, лучше не фиксировать
жёстко, а~каждый раз сверять с~первоисточником. Сюда относятся,
например:
\begin{itemize}
  \item актуальные пороги мониторинговых программ карточных сетей;
  \item ограничения на повторы и~даты вступления требований в~силу;
  \item актуальные доли методов оплаты и сравнительные таблицы;
  \item тарифы, пороги и~даты, зависящие от конкретного выпуска правил;
  \item полные таблицы кодов ответа (response codes) и~кодов причин (reason codes).
\end{itemize}
Этот принцип полезен и~книге, и~внутренним базам знаний.
Если материал меняется каждые полгода,
ему не место в~длинном повествовательном тексте без явного цикла
перепроверки. В~книге он остаётся как метод, архитектурное
разграничение или иллюстративный пример, либо уступает место прямой
ссылке на основной документ.

\section*{Как удерживать рабочий ритм}\label{sec:ps-tracking-rhythm}

Отслеживание изменений лучше выстраивать как ритуал, а~не как авральную реакцию на инцидент. Для большинства платёжных команд хватает
простого, но устойчивого ритма.

Раз в~неделю полезно просматривать входящие бюллетени,
регуляторные новости и обновления PCI SSC и EMVCo. Ежемесячно
нужно разбирать новые документы по схеме из этой главы
и обновлять карту рисков. Ежеквартально нужно пересматривать
справочные таблицы и~решать, что в~книге устарело, что
следует убрать или переписать в~следующей редакции и~какие места
текста стоит жёстче привязать к~первоисточнику. Перед крупным
релизом разумно перечитывать основные документы именно
по тем областям, которых этот релиз касается.

Такой ритм не делает команду безошибочной, но уменьшает
шанс узнать о~новом обязательном поле, пороге или ограничении задним числом~--- из всплеска отказов, штрафа или замечаний в~сертификационном отчёте.

\section*{Итоги}
\begin{itemize}
  \item Главная задача при работе с~источниками: найти нужный
    документ и определить его версию и~характер изменения.
  \item Бюллетени, заметки к~выпуску, обновления PCI SSC и EMVCo
    и~проекты регуляторных актов нужно читать как ранние
    сигналы к~действию, а~не складывать в~информационный шум.
  \item Если детали скрыты в~закрытом документе, сначала
    зафиксировать факт изменения и дату вступления в~силу,
    затем получить основной документ через участника сети
    и~только после этого приступать к~проектированию решения.
  \item Каждое существенное изменение полезно раскладывать
    по вопросам: к~какой плоскости оно относится, кого затрагивает,
    какое действие требует и~к~какому сроку.
  \item Книга должна держать инварианты и архитектуру, а~быстро
    стареющие пороги, цифры и~даты нужно перепроверять
    по первоисточникам, а~не превращать в~обещание вечной
    актуальности.
\end{itemize}
